\section{A sharp additive threshold for prime cubes}
Lee and O'Clarey's Theorem 1.3 states that Dirichlet GRH implies that
every integer $n>24$ is the sum of a prime and a positive integer
divisible by a prime square. The Lean result supplies their stated
hypothesis. What happens if the second summand must be divisible by a
prime \emph{cube}? An explicit least-prime bound gives a sharp
threshold, and a complete exception list.\hfill
\href{https://doi.org/10.1007/s00013-026-02275-6}{[Lee--O'Clarey, 2026]}

\paragraph{Theorem.} Every integer $n>10500$ has a representation
$n=p+kq^3$ with $p,q$ prime and $k\ge1$. Precisely 62 positive integers
have no such representation. They are
\[
\begin{gathered}
1,2,3,4,5,6,7,8,9,12,14,16,17,20,22,24,28,36,48,52,54,60,62,\\
72,76,78,96,102,108,120,126,150,174,180,216,240,270,294,312,342,\\
420,468,486,528,552,720,750,780,828,870,924,966,1080,1236,1260,\\
1500,1590,1800,1860,2520,8190,10500 .
\end{gathered}
\]
For example, $10501=5+1312\cdot 2^3$. The quantifier on $k$ matters:
requiring the second summand itself to be a cube gives a different problem.

\paragraph{Proof.}
Let $q$ be the least prime which does not divide $n$, and
$B_q=\prod_{\ell<q,\ \ell\text{ prime}}\ell$. Then $B_q\mid n$ and
$n$ is a unit modulo $q^3$. Write $P(a,m)$ for the least prime
congruent to $a\pmod m$. If $P(n,q^3)<n$, subtraction supplies the
required positive multiple. Lamzouri--Li--Soundararajan's Corollary 1.2
applies to every modulus $q^3$ and gives
\[
 P(n,q^3)\le [\varphi(q^3)\log(q^3)]^2<9q^8 .
\]
The estimate uses the Dirichlet GRH supplied by the repository.
\hfill\href{https://arxiv.org/pdf/1309.3595}{[Lamzouri--Li--Soundararajan]}

For every prime $q\ge43$, $B_q>9q^8$. Exact multiplication verifies
this for $43\le q\le257$. If the inequality holds for a prime $q\ge257$
and $r$ is the next prime, then $r<2q$ by Bertrand's postulate and
$B_r=qB_q>9q^9>9(2q)^8>9r^8$. This proves the claim by induction.
Thus every $n$ whose least missing prime is at least 43 has a representation.

For $q\le41$ set $M_q=\max_{(a,q)=1}P(a,q^3)$, with $a$ running
through all residue classes modulo $q^3$. A sieve produces the following
exact values by retaining the first prime in each reduced class.
\begin{center}
\begin{tabular}{r|rrrrrrr}
$q$ &2&3&5&7&11&13&17\\
$M_q$&17&109&1621&9371&47737&120401&349529\\ \hline
$q$&19&23&29&31&37&41&\\
$M_q$&540961&1494349&2367523&3812923&5897123&7802437&
\end{tabular}
\end{center}
An exception must satisfy $B_q\le n\le M_q$. For each of
$q=23,29,31,37,41$ the lower endpoint exceeds the upper. Thus it
suffices to decide representability for $1\le n\le540961$.
For every prime $q$ with $q^3+2\le540961$ and every residue $a\bmod q^3$,
find the first prime $r$ in that class with $r+q^3\le540961$, if one
exists. Mark $r+q^3,r+2q^3,\ldots$ up to the bound. The marked values
have explicit representations. Conversely, every representation occurs
among the marks, since its prime is at least the first prime in the
same class. This step includes the class of $q$ itself. The 62 unmarked
values are those displayed above. The accompanying Python program carries
out the sieve, the primorial checks and the final enumeration. It also
tests each listed exception a second time by subtracting every eligible
positive multiple of each prime cube. An independent bit-set sum of all primes
and all positive cube multiples agrees on every integer through $540961$.

The largest exception can also be checked by hand. For $n=10500$, the
eligible primes $q$ are $2,3,5,7,11,13,17,19$. The first four divide $n$,
so the only candidate prime summand would be $q$ itself, and in each case
$q^3\nmid n-q$. For the remaining four primes an even multiplier would
give an even summand, and $n-2=10498$ is divisible by none of their cubes.
For odd multipliers the possible positive differences, grouped by $q$,
are
\[
 \{9169,6507,3845,1183\},\quad \{8303,3909\},\quad
 \{5587\},\quad \{3641\}.
\]
They have prime divisors $53,3,5,7$, then $19,3$, then $37$ and $11$,
respectively. This gives a direct check of sharpness.

\paragraph{Recovering the square theorem.}
A multiple of $q^3$ is a multiple of $q^2$. Thus only the 62 cube
exceptions can fail to be a prime plus a positive multiple of a prime
square. Direct subtraction of prime-square multiples on these 62 inputs
gives exactly $1,2,3,4,5,8,24$ as square exceptions. This gives a second
proof of Lee and O'Clarey's explicit conclusion, using the same
least-prime estimate as the cube theorem.

For each fixed exponent $j\ge2$ the same argument reduces
$n=p+kq^j$ to finitely many checks. Indeed the least-prime bound is
at most $j^2 q^{2j+2}$, while $B_q$ eventually exceeds every fixed
power of $q$ (for example by the prime number theorem). This observation
does not require a claim of priority for qualitative finiteness. The
cube calculation identifies the last exception exactly.
